\documentclass[11pt, a4paper]{article}

% --- PAQUETES DE IDIOMA Y FORMATO ---
\usepackage[utf8]{inputenc}
\usepackage[spanish, es-tabla]{babel}
\usepackage[margin=2.5cm]{geometry}
\usepackage{fancyhdr}
\setlength{\headheight}{26pt}

\usepackage{amsmath, amsfonts, amssymb, amsthm}
\usepackage{booktabs}
\usepackage{array}
\usepackage{enumitem}
\usepackage{hyperref}
\usepackage{tikz}

% --- CONFIGURACIONES DE ENCABEZADO Y PIE DE PÁGINA ---
\pagestyle{fancy}
\fancyhf{}
% Encabezado
\fancyhead[L]{\textbf{Universidad San Francisco de Quito}\\Matemáticas Discretas}
% Pie de página
\fancyfoot[C]{Página \thepage}
\renewcommand{\headrulewidth}{0.4pt}
\renewcommand{\footrulewidth}{0.4pt}

% --- DATOS DEL DOCUMENTO ---
\title{\vspace{-1cm}\Huge\textbf{Reporte del proyecto: Algoritmos para el cálculo de todos los Testores Típicos de una matriz básica}\\[0.3cm]\Large Matemáticas Discretas}

\author{\textbf{Integrantes:} Martín Enríquez (00347250), Paula Gallardo (00343921), Eduardo Tapia (00346319)\\[0.1cm]
    \textbf{Profesor:} Eduardo Alba\\
    Universidad San Francisco de Quito (USFQ)}
\date{\today}

\begin{document}
    
    \maketitle
    \thispagestyle{fancy}
    
    \hrule
    \vspace{0.5cm}
    
    % ==============================================================================
    % SECCIÓN 1
    % ==============================================================================
        [\textbf{Ejercicio 1:}] [Enunciado del ejercicio]
        
        \textbf{Solución:} 
        
        \vspace{0.5cm}
        
        [\textbf{Ejercicio 2:}] [Enunciado del ejercicio]
        
        \textbf{Solución:} 
        
        \vspace{0.5cm}
    
    \vspace{0.5cm}
    \hrule
    \vspace{0.5cm}
    
    % ==============================================================================
    % SECCIÓN 2
    % ==============================================================================

    
\end{document}